\documentclass[11pt, a4paper]{article}
\usepackage[utf8]{inputenc}
\usepackage[T1]{fontenc}
\usepackage{amsmath, amssymb}
\usepackage{hyperref}
\usepackage{geometry}
\geometry{margin=1in}

% 定義 Pandoc 轉換需要的緊湊清單指令
\providecommand{\tightlist}{%
  \setlength{\itemsep}{0pt}\setlength{\parskip}{0pt}}

\title{Inference Creep: Decision Expansion and Governance Drift in AI-Assisted Software Engineering}
\author{Spark Tsai}
\date{February 2026}

\begin{document}

\maketitle

\noindent\textbf{Category:} cs.SE (Software Engineering)\\
\textbf{Keywords:} Inference Creep, AI-Assisted Software Engineering, Decision Accountability, Traceability, Probabilistic Systems, Decision Boundary, AI4SE, Human-AI Decision Making

\begin{abstract}
This paper introduces \textbf{Inference Creep} as an engineering-level governance phenomenon in AI-assisted software engineering: the gradual and unauthorized expansion of AI-mediated decision influence beyond its originally authorized boundary. Unlike traditional software risks that stem from faulty implementations or incorrect outputs, Inference Creep arises when AI systems participate in decisions without explicit authorization, traceability, or boundary definition. We argue that this risk cannot be adequately addressed by runtime controls or output filtering alone. Crucially, this paper demonstrates that Inference Creep is not merely a conceptual concern. It is \textbf{engineering-detectable, observable, and reproducible} through minimal structural criteria that link specifications and execution artifacts. By reframing AI risk as a \textbf{decision-stage governance problem}, this work positions Inference Creep as a distinct and falsifiable risk category within software engineering.
\end{abstract}

\section{Introduction}
\label{sec:introduction}
AI-assisted tools are now deeply embedded in modern software engineering workflows, including code generation, refactoring, testing, review, and deployment. While these systems are often framed as productivity aids, their influence increasingly extends into areas traditionally governed by human decision authority.

This shift introduces a structural risk that is insufficiently captured by existing discussions of model accuracy, bias, or hallucination. The core issue is not whether AI outputs are correct, but \textbf{whether AI is making or shaping decisions it was never authorized to influence}. 

This paper addresses that gap by defining \textbf{Inference Creep} as a software engineering phenomenon. It argues that current governance approaches focus disproportionately on runtime behavior, while neglecting the decision formation stage where governance boundaries should be enforced.

\section{From Feature Creep to Inference Creep}
In traditional software engineering, \textbf{Feature Creep} describes the uncontrolled expansion of system functionality beyond its original scope. Feature Creep is visible, measurable, and typically managed through requirement control and scope management.

\textbf{Inference Creep}, by contrast, does not introduce new features or capabilities. Instead, it alters \textbf{who participates in decisions and how far that participation extends}. Key distinctions include:

\begin{itemize}
\tightlist
\item No change in model architecture or parameters.
\item No explicit addition of system functionality.
\item A shift in decision influence without corresponding authorization.
\end{itemize}

Inference Creep therefore represents a \textbf{boundary violation}, not a functional expansion. The affected boundary is not an API or module interface, but a \textbf{decision boundary}.

\section{Decision Expansion and Boundary Erosion}
In AI-assisted development environments, decision influence often expands incrementally: suggestions become defaults, defaults become expectations, and expectations become de facto authority.

Because these transitions are gradual and probabilistic, they are rarely documented as explicit design changes. As a result, the system's decision boundary erodes without formal recognition. Unlike traditional software boundaries enforced through interfaces or contracts, decision boundaries in AI-assisted workflows often lack concrete artifacts, making violations difficult to detect using conventional engineering tools.

\section{Probabilistic Amplification and Risk Acceleration}
Large language models operate through probabilistic inference. When AI-generated suggestions influence subsequent decisions, early deviations can propagate and amplify across development stages. This creates \textbf{Risk Acceleration}, where a minor, unreviewed inference influences later design choices, and subsequent AI interactions build upon earlier decisions.

Importantly, this amplification is not a statistical anomaly but a structural property of probabilistic decision chains. Runtime correctness checks may catch individual failures, but they cannot prevent systemic drift originating at the decision formation stage.

\section{Engineering-Detectable Criteria for Inference Creep}
To avoid reducing Inference Creep to a purely semantic concept, we introduce a \textbf{minimal, engineering-level detection criterion} based on traceability between specification and execution.

\subsection{System Entities}
Let:
\begin{itemize}
\tightlist
\item $S_{\text{spec}}$ denote the set of authorized specifications or documented intents.
\item $E_{\text{exec}}$ denote the set of executed behaviors, code changes, or AI-mediated decisions.
\item Each element in both sets may optionally carry a \textbf{Trace ID}, serving as an anchor between intent and execution.
\end{itemize}

\subsection{Trace-Based Creep Criterion}
For any executed element $e \in E_{\text{exec}}$:
\begin{enumerate}
    \item \textbf{Authorized execution:} If $\exists s \in S_{\text{spec}}$ such that $\text{TraceID}(e) = \text{TraceID}(s)$.
    \item \textbf{Inference Creep:} If $\text{TraceID}(e) \notin \{ \text{TraceID}(s) \mid s \in S_{\text{spec}} \}$.
\end{enumerate}



\subsection{Engineering Interpretation}
This formulation yields immediate, practical interpretations:
\begin{itemize}
\tightlist
\item A specification exists, but no corresponding execution $\rightarrow$ \textit{Unimplemented intent}.
\item One specification maps to multiple executions $\rightarrow$ \textit{Decision expansion}.
\item Execution exists without specification $\rightarrow$ \textit{Unauthorized decision} (\textbf{Inference Creep}).
\end{itemize}

\section{Limits of Runtime Controls}
Runtime controls operate after decisions have already been formed. While necessary for operational safety, they cannot address inference creep because they react to outcomes rather than decision formation, and they filter content rather than authority. Inference creep is therefore a \textbf{decision-stage governance failure}.

\section{Governance Observability Gap and Orphaned Decisions}
A critical consequence of inference creep is the emergence of a \textbf{Governance Observability Gap}. In engineering terms, a decision that cannot be traced to an authorized source should be treated as an \textbf{orphaned decision}.

An orphaned decision is defined by \textbf{missing provenance}, not by incorrect output. A decision without traceability is not merely risky; it is ungoverned. From a software engineering perspective, orphaned decisions are analogous to dangling pointers: their existence signals structural integrity failure.

\section{Conclusion}
Inference Creep represents a distinct class of risk that originates at the decision formation stage. By providing an engineering-detectable definition grounded in traceability and authorization boundaries, this paper demonstrates that Inference Creep is a measurable governance failure. As AI systems increasingly participate in software development, the ability to detect unauthorized decision expansion will become a foundational requirement for trustworthy engineering practice.

\end{document}